% =============================================================================
% 12-pipelining-risc-cisc-notes.tex -- Lecture Notes for Week 12
% DERIVED ARTIFACT. Source of truth:
% Slides/CS401/12-pipelining-risc-cisc.tex.
% The Notes reorganize the deck by topic while preserving its content,
% notation, worked examples, and citations.
%
% Anchor reading:
%   * Patterson & Hennessy ARM 5e Ch. 1, §1.6 (p. 28); Ch. 4,
%     §§4.5--4.8 (pp. 283, 297, 316, 328); Ch. 6, §§6.2--6.3
%     (pp. 518, 523); App. D (pp. D-2--D-43).
%   * Stallings 11e Ch. 16 §16.4 (PDF pp. 667--685), Ch. 17 §17.5
%     (PDF p. 730) and §17.9 (PDF p. 752), and Ch. 20 §20.1 (PDF p. 847).
%     The repository keeps the historical Stallings2015 bibliography key.
%
% Compile from Notes/CS401 on Windows/MiKTeX:
%   TEXINPUTS="../../Preambles;" xelatex -interaction=nonstopmode 12-pipelining-risc-cisc-notes.tex
%   BIBINPUTS="../..;" bibtex 12-pipelining-risc-cisc-notes
%   TEXINPUTS="../../Preambles;" xelatex -interaction=nonstopmode 12-pipelining-risc-cisc-notes.tex
%   TEXINPUTS="../../Preambles;" xelatex -interaction=nonstopmode 12-pipelining-risc-cisc-notes.tex
% =============================================================================

\documentclass[11pt]{article}
\usepackage[margin=1in]{geometry}
\usepackage{array}
\usepackage{caption}
% =============================================================================
% Preambles/header.tex — shared LaTeX/Beamer preamble for this project
%
% Use from a Beamer lecture by prepending:
%     \input{header}
% and compiling with TEXINPUTS=../Preambles:$TEXINPUTS (the /compile-latex
% skill does this automatically).
%
% PALETTE CONTRACT. Color names defined here MUST match the names in
% Quarto/theme-template.scss so Beamer and Quarto renderings use the same
% palette. The scripts/check-palette-sync.sh script enforces this.
%
% To customize: edit the HEX values below AND the matching $variable in
% Quarto/theme-template.scss, then run scripts/check-palette-sync.sh.
% =============================================================================

% -----------------------------------------------------------------------------
% Engine detection + encoding
%
% Our /compile-latex skill uses XeLaTeX, where inputenc is unnecessary and
% can produce encoding warnings. Only load inputenc under pdfTeX.
% -----------------------------------------------------------------------------
\usepackage{iftex}
\ifPDFTeX
  \usepackage[utf8]{inputenc}
\fi

\usepackage{amsmath, amssymb, amsthm}
\usepackage{booktabs}
\usepackage{graphicx}

% -----------------------------------------------------------------------------
% PALETTE — mirrors Quarto/theme-template.scss variable names.
% Load xcolor and define colors BEFORE anything that references them
% (notably hyperref, hypersetup, and the Beamer color assignments).
% -----------------------------------------------------------------------------
\usepackage{xcolor}

\definecolor{primary-blue}{HTML}{012169}      % headings, accents
\definecolor{primary-gold}{HTML}{B9975B}      % emphasis, borders
\definecolor{highlight-yellow}{HTML}{F2A900}  % markers, alerts
\definecolor{light-bg}{HTML}{E8EDF5}          % title / section backgrounds
\definecolor{jet}{HTML}{1A1A1A}               % body text

% Semantic colors (used in diagrams and annotations)
\definecolor{positive}{HTML}{15803D}          % good / observed / identified
\definecolor{negative}{HTML}{B91C1C}          % bad / problematic / bias
\definecolor{neutral}{HTML}{525252}           % reference / context

% Highlight accents (match .hi-* classes in the SCSS)
\definecolor{hi-slate}{HTML}{314F4F}
\definecolor{hi-green}{HTML}{2E8B57}
\definecolor{hi-red}{HTML}{C41E3A}

% Semantic aliases so skill templates and snippets can write
% \textcolor{accent}{...} without hard-coding "primary-blue".
\colorlet{accent}{primary-blue}
\colorlet{accent2}{primary-gold}

% -----------------------------------------------------------------------------
% Hyperref — loaded AFTER xcolor + palette so link colors resolve.
% -----------------------------------------------------------------------------
\usepackage{hyperref}
\hypersetup{colorlinks=true, linkcolor=primary-blue, urlcolor=primary-blue,
            citecolor=primary-blue, pdfborder={0 0 0}}

% -----------------------------------------------------------------------------
% Course code — set once per deck (after \input{header}, before \begin{document})
% via \coursecode{CS401}. Shown in the footer on every slide so a deck's course
% is identifiable without opening the title page. Empty by default.
% -----------------------------------------------------------------------------
\makeatletter
\newcommand{\coursecode}[1]{\gdef\@coursecodetext{#1}}
\gdef\@coursecodetext{}
\makeatother

% -----------------------------------------------------------------------------
% Beamer theme assignments (only apply under Beamer).
%
% We detect Beamer via \@ifclassloaded{beamer}, which is the documented,
% non-internal way. This must be wrapped in \makeatletter/\makeatother
% because \@ifclassloaded contains an @-catcode character.
% -----------------------------------------------------------------------------
\makeatletter
\@ifclassloaded{beamer}{%
  \usefonttheme{professionalfonts}
  \setbeamercolor{structure}{fg=primary-blue}
  \setbeamercolor{frametitle}{fg=primary-blue}
  \setbeamercolor{title}{fg=primary-blue}
  \setbeamercolor{subtitle}{fg=primary-gold}
  \setbeamercolor{author}{fg=primary-blue}
  \setbeamercolor{institute}{fg=neutral}
  \setbeamercolor{date}{fg=primary-gold}
  \setbeamercolor{itemize item}{fg=primary-blue}
  \setbeamercolor{itemize subitem}{fg=primary-gold}
  \setbeamercolor{alerted text}{fg=primary-gold}
  \setbeamercolor{block title}{fg=white, bg=primary-blue}
  \setbeamercolor{block body}{bg=light-bg}
  % Semantic block variants: block=definition/concept (blue), exampleblock=
  % worked example (green/positive), alertblock=key takeaway or caution (gold).
  % Keep to <=2 colored boxes per slide across ALL three variants (INV-7).
  \setbeamercolor{block title example}{fg=white, bg=positive}
  \setbeamercolor{block body example}{bg=light-bg}
  \setbeamercolor{block title alerted}{fg=white, bg=primary-gold}
  \setbeamercolor{block body alerted}{bg=light-bg}

  % Minimal footer: course code (left) + page X / N (right), no navigation clutter
  \setbeamertemplate{navigation symbols}{}
  \setbeamertemplate{footline}{%
    \color{neutral}\footnotesize\hspace{0.7em}\@coursecodetext\hfill
    \insertframenumber/\inserttotalframenumber\hspace{0.7em}\vspace{0.3em}}
}{}
\makeatother

% -----------------------------------------------------------------------------
% TikZ setup — libraries the snippet gallery uses
% -----------------------------------------------------------------------------
\usepackage{tikz}
\usetikzlibrary{arrows.meta, positioning, calc, decorations.pathreplacing,
                fit, shapes.geometric, backgrounds}

% Shared TikZ styles — snippets and hand-written diagrams can pick these up
% via `\begin{tikzpicture}[every node/.style={...}]` or by naming specific
% styles (e.g., `\node[dag-node] (X) at (0,0) {X};`).
\tikzset{
  % Node styles for causal DAGs and similar
  dag-node/.style = {
    draw=primary-blue, fill=light-bg, rounded corners=2pt,
    minimum width=1.2cm, minimum height=0.9cm, align=center, font=\small
  },
  decision-node/.style = {
    draw=primary-gold, fill=white, diamond, aspect=1.8,
    minimum width=1.8cm, minimum height=1.2cm, align=center, font=\small,
    inner sep=1pt
  },
  % Edge styles — observed vs counterfactual
  observed-edge/.style = {-{Latex[length=2.5mm]}, thick, draw=jet},
  counterfactual-edge/.style = {-{Latex[length=2.5mm]}, thick, draw=neutral, dashed},
  confound-edge/.style = {-{Latex[length=2.5mm]}, thick, draw=negative, dashed},
  % Data-point styles for plots
  observed-dot/.style = {fill=primary-blue, draw=primary-blue, circle, inner sep=1.2pt},
  counterfactual-dot/.style = {fill=white, draw=neutral, circle, inner sep=1.2pt}
}

% -----------------------------------------------------------------------------
% Convenience macros
% -----------------------------------------------------------------------------
\newcommand{\muted}[1]{\textcolor{neutral}{#1}}
\newcommand{\key}[1]{\textcolor{primary-gold}{\textbf{#1}}}
\newcommand{\good}[1]{\textcolor{positive}{#1}}
\newcommand{\bad}[1]{\textcolor{negative}{#1}}

% Standout frame macro: full-bleed dark background for section transitions.
% Beamer-only (uses \begin{frame}/\setbeamercolor/\usebeamerfont) -- do not
% call from an article-class file (e.g. Notes/<CODE>/*-notes.tex). \muted,
% \key, \good, \bad, and \coursecode above are plain xcolor/gdef and are
% safe to use from either class; verified when Notes/article-class support
% was added.
% Expand: \transitionslide{Your transition title}
\newcommand{\transitionslide}[1]{%
  \begingroup
  \setbeamercolor{background canvas}{bg=primary-blue}
  \begin{frame}[plain]
    \centering
    \vfill
    {\usebeamerfont{title}\color{white}\Large #1\par}
    \vfill
  \end{frame}
  \endgroup
}

% Section divider frame (Beamer-only), an alternative to \transitionslide with a
% darker two-line style: near-black (jet) background, a small white label line
% above a larger gold title, and a thin gold rule along the bottom edge.
% Expand: \sectiondivider{Section 2}{Pointers}
\newcommand{\sectiondivider}[2]{%
  \begingroup
  \setbeamercolor{background canvas}{bg=jet}
  \begin{frame}[plain]
    \vspace*{3.0cm}
    {\color{white}\ttfamily\large #1\par}
    \vspace{0.5cm}
    {\color{primary-gold}\LARGE #2\par}
    \begin{tikzpicture}[remember picture, overlay]
      \draw[primary-gold, line width=2.5pt]
        ($(current page.south west)+(0,0.1)$) -- ($(current page.south east)+(0,0.1)$);
    \end{tikzpicture}
  \end{frame}
  \endgroup
}

% -----------------------------------------------------------------------------
% Channel watermark — "Concepts That Click" (YouTube).
%
% Article-class files (CompetitiveExam Q/A sets, Notes, Assignments) get a
% light diagonal draft watermark on every page plus a centered footer naming
% the channel. Beamer decks get a subtle TikZ background watermark instead
% (the footline already carries the course code). Centralized here so every
% MCQ PDF is branded without per-file edits.
% -----------------------------------------------------------------------------
\makeatletter
\@ifclassloaded{article}{%
  \usepackage{draftwatermark}
  \SetWatermarkText{Concepts That Click}
  \SetWatermarkScale{1}
  \SetWatermarkFontSize{2cm}
  \SetWatermarkLightness{0.86}
  \SetWatermarkAngle{45}
  \usepackage{fancyhdr}
  \pagestyle{fancy}
  \fancyhf{}
  \fancyfoot[C]{\footnotesize\textcolor{neutral}{Concepts That Click~|~YouTube: \href{https://www.youtube.com/@concepts.thatclick}{@concepts.thatclick}}}
  \renewcommand{\headrulewidth}{0pt}
  \renewcommand{\footrulewidth}{0pt}
  \fancypagestyle{plain}{%
    \fancyhf{}%
    \fancyfoot[C]{\footnotesize\textcolor{neutral}{Concepts That Click~|~YouTube: \href{https://www.youtube.com/@concepts.thatclick}{@concepts.thatclick}}}%
    \renewcommand{\headrulewidth}{0pt}%
    \renewcommand{\footrulewidth}{0pt}%
  }
}{}
\@ifclassloaded{beamer}{%
  \setbeamertemplate{background}{%
    \begin{tikzpicture}[remember picture, overlay]
      \node[rotate=45, opacity=0.055, align=center]
        at (current page.center)
        {\Huge\textcolor{neutral}{Concepts That Click}};
    \end{tikzpicture}%
  }
}{}
\makeatother 
\coursecode{CS401}

\renewcommand{\thesection}{12.\arabic{section}}
\renewcommand{\thefigure}{12.\arabic{figure}}
\newtheorem{example}{Example}[section]

\newenvironment{definitionbox}[1]{%
  \par\noindent\textbf{Definition (#1).}\ \itshape
}{\par\vspace{0.3em}}

\title{Instruction Pipelining and RISC vs. CISC\\\large Lecture Notes --- Week 12}
\author{Dr. Auqib Hamid Lone\\\small Assistant Professor in Computer Science \& Engineering}
\date{\today}

\begin{document}
\maketitle

\section{The Course Arc and the Performance Question}
\label{sec:arc}

Week 11 ended with a precise warning about overlap: parallel work is useful
only when the machine preserves the meaning of the instruction stream. A
five-stage pipeline can expose structural, data, and control hazards. The
usual remedies---forwarding, stalls, scheduling, prediction, and recovery---
trade hardware or control complexity for useful throughput. Week 12 widens the
question. Pipelining is the mechanism that moves instruction work through
overlapped stages; RISC and CISC are different design philosophies for
describing and implementing that work.

The first discipline is to measure a program rather than infer its speed from
one attractive label. For a workload with instruction count $\mathrm{IC}$,
average cycles per instruction $\mathrm{CPI}$, and clock rate $f$, the CPU time
is
\[
  T_{CPU}=\frac{\mathrm{IC}\times\mathrm{CPI}}{f}.
\]
This is the same performance equation used in Week 1 and in Patterson and
Hennessy, Ch.~1, §1.6 (p.~28)
\cite{PattersonHennessy2017_computer_organization_design}. A pipeline can
reduce the effective CPI of a long stream, but pipeline fill, drain, stalls,
flushes, cache misses, and the instruction count remain part of the measured
time. The instruction set, compiler, datapath, memory hierarchy, and workload
therefore meet in the final result; no single label wins in isolation.

For the motivating two-processor comparison, Processor A completes one
instruction per cycle at 2~GHz, so its CPI is 1 and its time per instruction is
$1/(2\times10^9)$ seconds. Processor B has a 3~GHz clock but an average CPI of
2, so its time per instruction is $2/(3\times10^9)$ seconds. Thus A takes
$0.5$~ns per instruction while B takes approximately $0.667$~ns: the higher
clock rate does not win because B spends more cycles per instruction.

The running comparison for this chapter is the operation
\texttt{A[i] = A[i] + 1}. The operation requires an element to be loaded, one
to be added, and the result to be stored. An ISA determines where those
operands live and how much of that work an architectural instruction names.
The chapter develops four connected questions: how pipeline timing determines
latency and throughput, what regularity or expressiveness buys an ISA, how a
workload determines the relevant trade-off, and how the whole course can be
recalled through a compact diagnostic procedure.

\section{Pipelining: Latency, Throughput, and Hazards}
\label{sec:pipeline}

\subsection{The five-stage model}

\begin{definitionbox}{Five-stage pipeline}
The teaching model divides an instruction path into IF (instruction fetch), ID
(decode and register read), EX (ALU operation or branch work), MEM (data memory
access), and WB (register write). The same instruction visits all five stages,
while different instructions occupy different stages simultaneously.
\end{definitionbox}

Pipeline registers separate the stages and carry an instruction's control and
data state across clock boundaries. The clock period is constrained by the
slowest stage plus the overhead of the pipeline registers. The five-stage
functions are therefore useful for timing reasoning even though an actual
processor may divide or combine them differently.

\begin{definitionbox}{Latency, throughput, and speedup}
\textbf{Latency} is the time from an instruction's IF stage to its WB stage.
\textbf{Throughput} is the interval between completed results in an instruction
stream. \textbf{Speedup} compares the total completion time of an overlapped
implementation with the total time of a non-overlapped baseline.
\end{definitionbox}

The distinction matters. Once a $k$-stage pipeline is full, an ideal stream
can complete one instruction every clock cycle, but one instruction still
traverses all $k$ stages. Thus the pipeline improves steady-state throughput;
it does not make the latency of every instruction one cycle.

\subsection{Fill, drain, and finite-stream speedup}

\begin{example}[Six independent instructions]
Assume $n=6$ independent instructions, $k=5$ balanced stages, one cycle per
stage, no hazards, and no extra pipeline-register cost. Without overlap, every
instruction occupies all five stage-cycles, so
\[
  T_{seq}=nk=6\times5=30\text{ cycles}.
\]

With overlap, the first instruction needs five cycles to pass through the
pipeline. Each of the remaining five instructions enters one cycle after its
predecessor, so the last instruction completes at
\[
  T_{pipe}=k+(n-1)=5+5=10\text{ cycles}.
\]
The complete schedule makes the fill and drain visible:
\begin{center}
\scriptsize
\begin{tabular}{@{}c|cccccccccc@{}}
\toprule
 & 1 & 2 & 3 & 4 & 5 & 6 & 7 & 8 & 9 & 10\\
\midrule
$I_1$ & IF & ID & EX & MEM & WB & & & & &\\
$I_2$ & & IF & ID & EX & MEM & WB & & & &\\
$I_3$ & & & IF & ID & EX & MEM & WB & & &\\
$I_4$ & & & & IF & ID & EX & MEM & WB & &\\
$I_5$ & & & & & IF & ID & EX & MEM & WB &\\
$I_6$ & & & & & & IF & ID & EX & MEM & WB\\
\bottomrule
\end{tabular}
\end{center}
The finite-stream speedup is therefore
\[
  S=\frac{30}{10}=3.0\times.
\]
The result is below the ideal fivefold speedup because this finite stream pays
for filling the stages at the beginning and draining them at the end.
\end{example}

For a general independent stream, the same reasoning gives
\[
  T_{pipe}=\bigl(k+n-1\bigr)T_c,
  \qquad
  S_{ideal}=\frac{nk}{k+n-1},
\]
where $T_c$ is the clock period. The ideal latency is approximately $kT_c$,
whereas the ideal steady-state throughput is one completed instruction per
$T_c$. As $n\to\infty$, $S_{ideal}\to k$ only when the stages are balanced and
the stream supplies enough independent work.

\subsection{Bubbles and the load-use dependence}

\begin{example}[A pipeline with one bubble]
Assume that a load's data becomes available after MEM. The immediately
following ALU instruction needs that value at EX, so forwarding alone cannot
make the value arrive soon enough. Consider
\[
\texttt{LW R1,0(R2)},\quad
\texttt{ADD R3,R1,R4},\quad
\texttt{SUB R5,R3,R6},\quad
\texttt{OR R7,R8,R9}.
\]
The schedule is:
\begin{center}
\scriptsize
\begin{tabular}{@{}c|ccccccccc@{}}
\toprule
 & 1 & 2 & 3 & 4 & 5 & 6 & 7 & 8 & 9\\
\midrule
\texttt{LW R1,0(R2)} & IF & ID & EX & MEM & WB & & & &\\
\texttt{ADD R3,R1,R4} & & IF & ID & \textbf{stall} & EX & MEM & WB & &\\
\texttt{SUB R5,R3,R6} & & & IF & \textbf{stall} & ID & EX & MEM & WB &\\
\texttt{OR R7,R8,R9} & & & & & IF & ID & EX & MEM & WB\\
\bottomrule
\end{tabular}
\end{center}

Without the bubble, the ADD would request the value before MEM had produced
it. The dependence is a true read-after-write (RAW) hazard. The one bubble
delays the ADD until the load result can be forwarded to its EX input; the
stream therefore completes in 9 rather than the ideal 8 cycles. Its achieved
CPI is shaped by the instruction mix and hazard rate, not only by the nominal
number of stages.
\end{example}

The Socratic conclusion is consequently precise: a five-stage pipeline that
completes six instructions in ten cycles has not made each instruction five
times faster. The total time fell from 30 to 10 cycles, giving a $3\times$
speedup for this finite workload. A complete performance report separates
latency and throughput, then accounts for bubbles, flushes, memory stalls, and
fill/drain effects.

\subsection{Pipeline design knobs}

Each design choice exchanges a possible gain for a new cost.
\begin{center}
\small
\begin{tabular}{@{}>{\raggedright\arraybackslash}p{2.8cm}>{\raggedright\arraybackslash}p{3.3cm}>{\raggedright\arraybackslash}p{3.3cm}>{\raggedright\arraybackslash}p{2.8cm}@{}}
\toprule
\textbf{Choice} & \textbf{Potential gain} & \textbf{New cost} & \textbf{Question}\\
\midrule
Deeper pipeline & Shorter stage delay; higher $f$ & More registers and branch penalty & Is useful work predictable?\\
Duplicate resources & Fewer structural hazards & Area, power, verification & Is the conflict frequent?\\
Forwarding & Fewer RAW stalls & Muxes, comparators, timing paths & Does the value arrive in time?\\
Prediction & Fewer control bubbles & Predictor state and flush recovery & What is $P_{flush}$?\\
\bottomrule
\end{tabular}
\end{center}
Stallings 11e treats instruction pipelining in §16.4 (PDF pp.~667--685); its
Ch.~17 RISC discussion connects regular instruction execution to pipelining in
§17.5 (PDF p.~730). These are indexed chapter and section anchors rather than
claims that the constructed examples are quotations from the books.

\section{RISC and CISC as Design Philosophies}
\label{sec:risc-cisc}

\begin{definitionbox}{RISC and CISC}
\textbf{RISC} is a tendency toward a relatively small, regular instruction set
with simple operations and a register-centered execution model. \textbf{CISC}
is a tendency toward a richer instruction set in which an individual
instruction may express more complex operations or addressing forms.
\end{definitionbox}

These are tendencies, not universal rankings. It is too strong to say that
every RISC instruction is simple or every CISC instruction is slow. A fair
comparison concerns the complete implementation: encoding, decode,
micro-operations, compiler, memory system, and target workload. Stallings
11e discusses RISC execution, RISC pipelining, and the RISC/CISC comparison in
§17.9 (PDF p.~752).

\subsection{The running operation in two instruction streams}

For \texttt{A[i] = A[i] + 1}, a register-centered or RISC tendency might make
the load, add, and store explicit:
\begin{center}
\texttt{LOAD R1,[R2+R3]}\\
\texttt{ADD R1,R1,1}\\
\texttt{STORE [R2+R3],R1}
\end{center}
A memory-rich or CISC tendency might expose a compact architectural operation
such as \texttt{INC [R2+R3]}. The operation is the same, but the visible
instruction stream differs. The CISC-like instruction may be split internally
into simpler micro-operations before execution. This is a conceptual
comparison, not a claim about one specific ISA's exact encoding. Patterson and
Hennessy App.~D surveys RISC addressing modes and formats (pp.~D-2--D-43).
The comparison can be tabulated as follows:
\begin{center}
\small
\begin{tabular}{@{}>{\raggedright\arraybackslash}p{5.6cm}>{\raggedright\arraybackslash}p{5.6cm}@{}}
\toprule
\textbf{Register-centered / RISC tendency} & \textbf{Memory-rich / CISC tendency}\\
\midrule
\texttt{LOAD R1,[R2+R3]} & \texttt{INC [R2+R3]}\\
\texttt{ADD R1,R1,1} & \textit{one architectural instruction may name the memory update}\\
\texttt{STORE [R2+R3],R1} & \textit{decode/execute may internally split the work}\\
More instructions, regular formats, explicit loads/stores & Fewer visible instructions, richer formats, more decode work\\
\bottomrule
\end{tabular}
\end{center}

Regularity can help a pipeline in several related ways:
\begin{itemize}
  \item Fixed or clustered instruction formats simplify fetch boundaries and
        decode timing.
  \item A large register file reduces frequent memory operands and exposes
        values to forwarding paths. This is a RISC design tendency, not a law.
  \item Load/store separation makes the memory stage's role predictable: ALU
        instructions can use EX without also requesting data memory.
  \item Predictable stage work makes balancing and hazard detection easier,
        while the compiler carries more scheduling responsibility.
\end{itemize}
The cost can move into the rest of the system. More instructions can increase
$\mathrm{IC}$, register traffic, code size, instruction-cache pressure, or
branch exposure for a particular workload.

Expressiveness can also be valuable. A complex instruction may reduce the
visible instruction count for a dense operation or a legacy software
interface. Rich addressing modes can express array, pointer, stack, or string
operations compactly. Compatibility and code density matter when instruction
memory, binary size, or an established software ecosystem dominates. The cost
does not disappear: decode width, variable-length boundaries, micro-operation
expansion, precise exceptions, and verification can complicate the front end.

\begin{example}[Avoiding an overstrong comparison]
The following statements are reasonable tendencies:
\begin{enumerate}
  \item RISC often favors regular formats and register operands.
  \item CISC can express more work per architectural instruction.
\end{enumerate}
The statement ``RISC always has fewer instructions and CISC always has lower
CPI'' is too strong. Instruction count, CPI, clock rate, stalls, code size,
and workload interact through $T_{CPU}=\mathrm{IC}\times\mathrm{CPI}/f$.
The appropriate comparison fixes the workload and measures total execution
time rather than selecting one favorable ISA statistic.
\end{example}

\section{Workload-Based Architectural Choice}
\label{sec:workload}

\begin{example}[Two workload sketches]
Consider an embedded-control workload with a tight power budget, predictable
loops, limited memory, and fixed sensor operations. Its likely priorities are
predictability, energy, and a compact implementation. A regular pipeline,
simple decode, and compiler-visible registers may help, while code expansion
and instruction-memory pressure are risks.

Now consider cloud analytics with a large software base, diverse branches,
high throughput, and frequent vector or multicore parallelism. Compatibility,
scalable parallel work, and a capable front end may matter more. Decode and
recovery complexity and memory bottlenecks remain risks. Patterson and
Hennessy Ch.~6 emphasizes that parallel performance depends on the program and
its communication and coordination costs (pp.~516--579).

The two sketches can be compared directly:
\begin{center}
\scriptsize
\begin{tabular}{@{}>{\raggedright\arraybackslash}p{3.0cm}>{\raggedright\arraybackslash}p{4.0cm}>{\raggedright\arraybackslash}p{4.0cm}@{}}
\toprule
\textbf{Criterion} & \textbf{Embedded control} & \textbf{Cloud analytics}\\
\midrule
Likely priority & Predictability, energy, compact implementation & Throughput, compatibility, scalable parallel work\\
Helpful tendency & Regular pipeline, simple decode, compiler-visible registers & Strong ecosystem, capable front end, SIMD/vector and multicore support\\
Main risk & Code expansion and instruction-memory pressure & Decode/recovery complexity and memory bottlenecks\\
\bottomrule
\end{tabular}
\end{center}

Neither row proves a universal winner. Each row identifies what should be
measured for the workload.
\end{example}

\begin{example}[Compare total time, not labels]
A RISC-like implementation executes $1.20\times10^9$ instructions at
$\mathrm{CPI}=1.1$ and $f=2.4$~GHz. A CISC-like implementation executes
$0.80\times10^9$ instructions at $\mathrm{CPI}=2.0$ and $f=3.0$~GHz. Applying
the performance equation gives
\begin{align*}
  T_R &= \frac{1.20\times10^9\times1.1}{2.4\times10^9}
       =0.55~\text{s},\\
  T_C &= \frac{0.80\times10^9\times2.0}{3.0\times10^9}
       \approx0.533~\text{s}.
\end{align*}
The CISC-like system wins this workload by about $3.0\%$ despite its higher
CPI, because it executes fewer instructions at a higher clock. A different
workload, cache behavior, or branch rate can reverse the result. The
measurement is the conclusion; the labels are hypotheses.
\end{example}

\subsection{The ISA boundary and the pipeline}

ISA decisions reappear as pipeline constraints.
\begin{center}
\small
\begin{tabular}{@{}>{\raggedright\arraybackslash}p{3.1cm}>{\raggedright\arraybackslash}p{4.2cm}>{\raggedright\arraybackslash}p{4.2cm}@{}}
\toprule
\textbf{Pipeline concern} & \textbf{Regular instruction stream helps} & \textbf{Rich instruction stream may require}\\
\midrule
Fetch/decode & Predictable boundaries and widths & Boundary detection and wider decode\\
Data hazards & Explicit register dependencies & Internal micro-op dependencies\\
Structural hazards & Regular stage/resource demand & Variable resource demand per instruction\\
Control recovery & Explicit branch behavior & Precise recovery across expanded work\\
\bottomrule
\end{tabular}
\end{center}

This is the course-level connection. Week 5's addressing modes, Weeks 6--7's
control units, Weeks 9--10's memory hierarchy, and Week 11's hazards all
reappear as constraints on instruction throughput. The final design question
is always ``fit to which workload?''

\section{Course Review and Exam Procedure}
\label{sec:review}

\subsection{One performance thread through CS401}

The course can be read as a sequence of increasingly complete performance
questions:
\begin{enumerate}
  \item \key{Representation:} bits, two's complement, and IEEE 754 determine
        what values hardware can encode.
  \item \key{Execution:} registers, buses, addressing modes, and control units
        move and transform those values.
  \item \key{Communication:} I/O and interrupts connect the processor to
        devices without making every transfer a busy wait.
  \item \key{Storage:} caches, virtual memory, and secondary storage trade
        capacity, latency, protection, and persistence.
  \item \key{Overlap:} parallel organizations and pipelines raise throughput
        when hazards are identified and paid for explicitly.
  \item \key{Choice:} the ISA and organization should be judged by the
        workload's measured time, energy, reliability, and software constraints.
\end{enumerate}

When a problem looks unfamiliar, ask six questions:
\begin{enumerate}
  \item What is the representation: bits, fields, signedness, or format?
  \item Where is the operand: register, memory, cache, page, or device?
  \item Which hardware state changes: PC, IR, MAR/MDR, register, flag, or
        control word?
  \item What is the bottleneck: clock, CPI, bus, memory, I/O, or hazard?
  \item What is the unit of comparison: latency, throughput, speedup, hit
        rate, miss penalty, or total CPU time?
  \item Which assumption makes the shortcut valid, and what breaks it?
\end{enumerate}
A strong exam solution names the model before doing arithmetic.

\begin{example}[Integrating the course topics]
Suppose a workload has many independent vector operations, frequent branches,
and a large working set. It runs on a multicore processor with a pipelined
front end. Inspect Week 11 and Week 12 for SIMD or MIMD opportunity, pipeline
hazards, branch penalty, and instruction-stream choice. Inspect Weeks 9--10
for cache locality, TLB/page behavior, and storage misses that may dominate an
otherwise fast pipeline. Finally, return to Week 1 and evaluate total time
through $\mathrm{IC}\times\mathrm{CPI}/f$ rather than inferring performance from
clock rate alone.
\end{example}

For exam and assignment problems, the same procedure becomes a checklist:
\begin{itemize}
  \item \textbf{Performance:} write the equation, normalize units, and
        compare total time or speedup.
  \item \textbf{Instruction cycle/control:} name the state and write RTL
        transfers with \texttt{PC}, \texttt{IR}, \texttt{MAR}, and \texttt{MDR}.
  \item \textbf{Memory:} split addresses into tag/index/offset or page/offset
        before translating or computing a hit.
  \item \textbf{I/O:} identify who waits, who owns the bus, and where the
        data moves.
  \item \textbf{Pipeline:} draw a cycle table; classify each hazard; count
        every stall and flush.
  \item \textbf{RISC/CISC:} state the workload, measured variables, and
        trade-off before defending a design choice.
\end{itemize}

The recurring course question is how a finite amount of hardware can keep
common work moving quickly, correctly, and predictably. Representation makes
operations possible; organization moves operands and generates control;
hierarchy hides slow capacity behind fast common cases; parallelism overlaps
work while hazards expose dependencies; and RISC/CISC choices shape the
instruction work. The workload decides whether the trade-off pays. Therefore,
do not ask which architecture is universally best. Ask which design makes
this workload's useful work per unit time highest, and state the assumptions
that make the answer true.

\subsection*{Exercises}
\begin{enumerate}
  \item Six independent instructions run on a balanced five-stage pipeline,
        one cycle per stage. Compute sequential time, ideal pipelined time,
        finite-stream speedup, and the limiting large-stream speedup.
  \item For the sequence \texttt{LW R1,0(R2)}, \texttt{ADD R3,R1,R4},
        \texttt{SUB R5,R3,R6}, and \texttt{OR R7,R8,R9}, identify the true
        dependences, insert the required bubble, and state which result can be
        forwarded.
  \item Explain why ``RISC always has fewer instructions and CISC always has
        lower CPI'' is not a valid universal conclusion. Name the variables
        that must be measured instead.
  \item A RISC-like processor executes $0.90\times10^9$ instructions at
        CPI $=1.4$ and $f=2.1$~GHz. A CISC-like processor executes
        $0.60\times10^9$ instructions at CPI $=2.2$ and $f=2.8$~GHz. Compute
        both CPU times and identify the faster implementation for this
        workload.
\end{enumerate}

\subsection*{Solutions}
\begingroup\small
\begin{enumerate}
  \item Sequential time is $6\times5=30$ cycles. Ideal pipelined time is
        $5+(6-1)=10$ cycles. Speedup is $30/10=3.0\times$. As the stream
        becomes very long, the ideal speedup approaches five, provided the
        stages are balanced and the stream supplies independent work.
  \item The load produces $R1$, which ADD reads, so the first edge is a true
        RAW load-use dependence and requires one bubble before ADD's EX stage.
        ADD produces $R3$, which SUB reads; that ALU result can be forwarded
        to SUB without an additional bubble. The final OR is independent of
        those values in the shown stream.
  \item The performance equation is
        $T_{CPU}=\mathrm{IC}\times\mathrm{CPI}/f$. Instruction count, CPI,
        clock rate, stalls, code size, cache behavior, branch behavior, and
        the workload interact, so neither ISA label determines total time by
        itself.
  \item For the RISC-like processor,
        \[
          T_R=\frac{0.90\times10^9\times1.4}{2.1\times10^9}=0.60~\text{s}.
        \]
        For the CISC-like processor,
        \[
          T_C=\frac{0.60\times10^9\times2.2}{2.8\times10^9}=0.4714~\text{s}
          \text{ (approximately)}.
        \]
        The CISC-like implementation is faster for this workload because its
        lower instruction count and higher clock outweigh its higher CPI.
\end{enumerate}
\endgroup

\section*{References and Source Boundaries}

The deck's references are retained without changing their scope. Stallings,
\emph{Computer Organization and Architecture: Designing for Performance},
11th ed. (the repository key is retained as Stallings2015), is used for
instruction pipelining in Ch.~16 §16.4 (PDF pp.~667--685), RISC pipelining in
Ch.~17 §17.5 (PDF p.~730), the RISC/CISC comparison in §17.9 (PDF p.~752),
and the brief stream-taxonomy connection in Ch.~20 §20.1 (PDF p.~847).
Patterson and Hennessy, \emph{Computer Organization and Design: The
Hardware/Software Interface}, ARM 5e, is used for the performance equation in
Ch.~1 §1.6 (p.~28), the five-stage and hazard model in Ch.~4 §§4.5--4.8
(pp.~283, 297, 316, 328), workload/parallelism framing in Ch.~6 §§6.2--6.3
(pp.~518, 523), and the RISC survey in App.~D (pp.~D-2--D-43). The worked
comparisons in this chapter are constructed examples: their numbers illustrate
the performance equation and pipeline model, and are not quotations or
measurements from either book.

\begingroup
\raggedright
\bibliographystyle{plain}
\bibliography{../../Bibliography_base}
\endgroup

\end{document}